\section{Related Work}
\emph{Security foundations.} Least privilege and complete mediation~\cite{saltzer}, privilege separation~\cite{privsep}, and capability-based sandboxing~\cite{capsicum} establish ways to limit authority and isolate execution. Models of information flow~\cite{denning} govern data propagation, while enforceable-policy theory~\cite{schneider} characterizes runtime monitoring guarantees. Together, these foundations connect permission assignment with control over actions and information.

\noindent\emph{Agent authorization and user control.} Surveys and vision papers examine user-level permission specification and enforcement~\cite{permissions}, authorization architectures~\cite{architectures}, and contextual information release~\cite{aac}. Complementary studies investigate contextual permission preferences~\cite{dataperms}, user involvement through Janus~\cite{janus}, and reusable user-authored policies~\cite{userpolicies}. Prudentia~\cite{prudentia} uses policy-aware planning to reduce approval requests, while elastic least privilege~\cite{elasticprivilege} proposes making authority costs visible to planners and limiting temporary grants.

\noindent\emph{Agent security evaluation.} AgentDojo~\cite{agentdojo} evaluates prompt-injection attacks and defenses in stateful tasks, separating task utility from attack outcomes. ToolEmu~\cite{toolemu} emulates tool execution with language models to identify potential failures. ToolPrivBench~\cite{toolprivbench} evaluates whether agents choose higher-privilege tools despite sufficient lower-privilege alternatives. These benchmarks provide foundations for evaluating authorization across alternative plans, changing grants, and external effects.

\noindent\emph{Requirements and software ecosystems.} Harel and Marron~\cite{requirements} examine the difficulty of discovering and describing intended behavior, beyond implementing a given specification. This perspective helps separate two authorization questions: whether a policy expresses the user's intended delegation and whether execution respects that policy. Zhao et al.~\cite{appstore} study LLM app stores through stakeholder relationships, security risks, and governance needs. Their ecosystem perspective is relevant when agents combine services supplied by different developers and platforms. Together, these perspectives connect the content of authorization requirements with responsibility for defining and enforcing them. They motivate explicit treatment of whose decisions a policy represents and where those decisions take effect across service boundaries.
